% 每个用例在 \documentclass 之后先读这个文件。
%
% 类在文末把没排过的前置、正文、后置补上（src/engine/hit-structure.dtx）。用例只测一样东西，
% 不要这些：三个标志位先置上，要哪段就在用例里自己写 \hitfrontmatter 那几句。
\ExplSyntaxOn
\bool_gset_true:N \g__hit_matter_front_bool
\bool_gset_true:N \g__hit_matter_main_bool
\bool_gset_true:N \g__hit_matter_back_bool
%
% 几个由类按分支定的开关写进日志，检查脚本用 t:fact(名字) 取（shenzhen 是改排之后的：
% 深圳本科终稿走本部模板，这里是假）。用例测的是按开关排得对不对，
% 开关各分支该取什么另有用例管，检查脚本里不再把那套规则抄一遍。文末才写，取的是排完以后的值。
\cs_new_protected:Npn \__feature_fact:nN #1#2
  { \iow_log:e { FEATURE-FACT~#1=\bool_if:NTF #2 { true } { false } } }
\AtEndDocument
  {
    \__feature_fact:nN { chapter } \g__hit_chapter_bool
    \__feature_fact:nN { en } \g__hit_en_bool
    \__feature_fact:nN { shenzhen } \g__hit_shenzhen_bool
    \__feature_fact:nN { bilingual-caption } \g__hit_bilingual_caption_bool
  }
\ExplSyntaxOff
